\ifdefined\iscombined
	\lecturetitle{Testing Levels}
\else
\documentclass{article}
\usepackage{listings}
\usepackage{amsthm}
\usepackage{comment}
\usepackage{amsmath}
\usepackage{amsfonts}
\usepackage{sectsty}
\usepackage{hyperref}
\usepackage{fancyhdr}
\usepackage[top=1in,bottom=1in,left=1in,right=1in]{geometry}
\usepackage{float}
\usepackage{graphicx}
\usepackage{tabularx}
\usepackage{mathtools}

\date{
	\vspace{-.75em}
	{\large \today}
}

\title{
	\vspace*{-1.5em}
	{\Huge Lecture 4} \\
	\vspace{-1em}
	\rule{\textwidth/4}{.01em} \\
	\vspace{-.25em}
	{\Large Testing Levels}
	\vspace{-.75em}
}

\author{
	{\large Matthew Farah}
}

\hypersetup{
	colorlinks=true,
	linkcolor=blue,
	filecolor=magenta,
	urlcolor=blue,
	citecolor=blue,
	anchorcolor=blue
}

\fancyhead{}
\fancyhead[L]{\today}
\fancyhead[R]{Lecture 4 - Testing Levels}

\sectionfont{\fontsize{12}{12}\bfseries}
\subsectionfont{\fontsize{10}{12}\normalfont}
\subsubsectionfont{\fontsize{10}{12}\normalfont}

\begin{document}

\maketitle
\pagestyle{fancy}

\fi

\begin{itemize}
	\item A test's `testing level' describe the scope of the software being tested.
	\item The three testing levels used in the course are:
	\begin{enumerate}
		\item Unit.
		\begin{itemize}
			\item Tests small, meaningful pieces of the code.
			\item Always done as white-box tests.
			\item Generally inexpensive to develop.
			\item Should be run as often as possible.
		\end{itemize}
		\item Integration.
		\begin{itemize}
			\item Tests how units work together.
			\item Done by treating units are black-boxes.
			\item Low to moderately expense to develop.
			\item Should be done after the creation of every new unit.
		\end{itemize}
		\item System.
		\begin{itemize}
			\item Tests the entire system.
			\item May be done as white or black box tests.
			\item Important functional and non-functional tests are typically done at this level.
			\item Most expensive testing level.
			\item Should be done prior to every release.
		\end{itemize}
	\end{enumerate}
	\item There are three primary approaches to designing good integration tests.
	\begin{enumerate}
		\item The big-bang approach:
		\begin{itemize}
			\item Create integration tests for all units at once.
			\item TLDR; shit, dumb, cringe.
		\end{itemize}
		\item The top-down approach:
		\begin{itemize}
			\item Create integration tests starting with the largest units/components first.
			\item TLDR; Mid asf.
		\end{itemize}
		\item The bottom-up approach:
		\begin{itemize}
			\item Create integration tests starting with the smallest units/components first.
			\item TLDR; based, epic.
		\end{itemize}
	\end{enumerate}
	\item System tests are designed to do one of two things:
	\begin{enumerate}
		\item Verify the system.
		\begin{itemize}
			\item Ensure that the system meets agreed upon specifications.
		\end{itemize}
		\item Validate the system.
		\begin{itemize}
			\item Ensure that the system delivers on what its user(s) want.
			\item Validating system tests are also called acceptance tests.
		\end{itemize}
	\end{enumerate}
	\item Tests can be developed manually or by an automated process.
	\begin{itemize}
		\item Context and the kind of test being developed dictate the method by which it should be developed.
	\end{itemize}
\end{itemize}

\ifdefined\iscombined
	\newpage
\else
	\end{document}
\fi
